%% ============================================================
%% AULA 11 — Memória e mapas: Hopfield, Kohonen e aprendizado
%%           hebbiano
%% GBC073 — Inteligência Computacional (FACOM/UFU)
%%
%% Esqueleto com esboço de conteúdo — a desenvolver.
%% Ementa (ficha SEI 5116656): item 2 — arquiteturas Hopfield e
%% Kohonen; aprendizagem não-supervisionada;
%% aplicações (clusterização).
%%
%% Compilar:  latexmk -xelatex aula11.tex   (duas passadas!)
%% ============================================================
\documentclass[aspectratio=169, 11pt]{beamer}

\makeatletter
\def\input@path{{./}{../}}
\makeatother

\usetheme{InteligenciaComputacional}   % ANTES do babel: fontspec vem primeiro
\usepackage{babel}
\babelprovide[import, main]{brazilian}

\title[Aula 11 — Hopfield e Kohonen]{Aula 11: Memória e mapas — Hopfield, Kohonen e aprendizado hebbiano}
\subtitle[GBC073 — Inteligência Computacional]{Hebb \textbullet\ memória associativa \textbullet\ energia \textbullet\ mapas auto-organizáveis}
\author[Prof.\ Marcelo Albertini]{Prof.\ Marcelo Keese Albertini}
\institute{GBC073 — Inteligência Computacional \textbullet\ Universidade Federal de Uberlândia}
\date{2026}

\begin{document}

% ------------------------------------------------------------ capa
\begin{frame}[plain, noframenumbering]
  \titlepage
\end{frame}

% ------------------------------------------------------------ roteiro
\begin{frame}{Roteiro da aula}
  \begin{itemize}\setlength{\itemsep}{0.5ex}
    \item<1-> \textbf{A regra de Hebb:} neurônios que disparam juntos, se conectam
    \item<2-> \textbf{Hopfield:} memória associativa e função de energia
    \item<3-> \textbf{Kohonen:} mapas auto-organizáveis e aprendizado competitivo
    \item<4-> \textbf{Hopfield moderno:} a ponte para a atenção
  \end{itemize}
\end{frame}

% ------------------------------------------------------------
\section{A regra de Hebb}

\begin{frame}{A regra de Hebb (neurônios que disparam juntos)}
  \begin{itemize}\setlength{\itemsep}{0.4ex}
    \item ``Neurônios que disparam juntos, se conectam'': Hebb (1949).
    \item Regra de aprendizado local: \(\Delta \pesoij{j}{i} = \taxa\, y_j\, y_i\)
          — sem rótulo, sem erro.
    \item É aprendizado por \emph{correlação}, não por correção.
    \item Raiz comum do que vem a seguir: Hopfield (memória) e Kohonen
          (mapas).
  \end{itemize}
\end{frame}

% ------------------------------------------------------------
\section{Hopfield: memória associativa e energia}

\begin{frame}{Hopfield: memória associativa e energia}
  \begin{itemize}\setlength{\itemsep}{0.4ex}
    \item Memória associativa: recordar o padrão inteiro a partir de um
          pedaço (ou de uma versão corrompida).
    \item Pesos simétricos, atualização assíncrona — a dinâmica desce a
          função de energia.
    \item Padrões memorizados viram mínimos de energia; a rede ``cai'' no
          mais próximo.
    \item Capacidade limitada e estados espúrios: o preço da simplicidade.
  \end{itemize}
  \begin{alertabox}
    A rede de Hopfield não ``raciocina'': ela desce uma superfície de energia
    até o padrão memorizado mais próximo — nada mais, nada menos.
  \end{alertabox}
\end{frame}

% ------------------------------------------------------------
\section{Kohonen: mapas auto-organizáveis}

\begin{frame}{Kohonen: mapas auto-organizáveis e aprendizado competitivo}
  \begin{itemize}\setlength{\itemsep}{0.4ex}
    \item Mapa auto-organizável (SOM): uma grade de neurônios que se molda
          aos dados.
    \item Aprendizado competitivo: o neurônio mais próximo do exemplo vence —
          e atualiza a si e aos vizinhos.
    \item A grade preserva a topologia dos dados: entradas parecidas ativam
          regiões vizinhas.
    \item Uso clássico: visualização e clusterização de dados de alta
          dimensão.
  \end{itemize}
  \begin{ideiabox}
    O SOM é clusterização com uma restrição a mais: os clusters precisam
    morar em uma grade — o que a torna também um mapa.
  \end{ideiabox}
\end{frame}

% ------------------------------------------------------------
\section{Hopfield moderno}

\begin{frame}{Hopfield moderno: a ponte para a atenção}
  \begin{itemize}\setlength{\itemsep}{0.4ex}
    \item Hopfield moderno (Ramsauer et al., 2020): atualização exponencial,
          memória contínua.
    \item Equivalência com a atenção: o self-attention é uma rede de Hopfield
          moderna.
    \item As memórias dos anos 1980 reaparecem dentro dos transformers.
    \item Fechando o ciclo com a aula 9 — e com a própria ficha da
          disciplina.
  \end{itemize}
\end{frame}

\end{document}
